\section{프로젝트 개요}

\subsection{프로젝트 정의와 목표}

ARES(Autonomous Rover Exploration System)는 초등학생을 주 대상으로 하는
교육용 로봇 코딩 플랫폼이다. 학습자는 웹 브라우저에서 Google Blockly\footnote{\textbf{Blockly}:
Google이 개발한 오픈소스 시각적 프로그래밍 라이브러리. 퍼즐 형태의 블록을 끼워 맞춰
프로그램을 구성하며, 조립된 블록은 실행 가능한 코드나 명령으로 변환된다.} 기반의
시각적 블록을 조립하여 프로그램을 작성하고, 이를 Bluetooth로 라즈베리파이
피코(Raspberry Pi Pico) 기반 로버 ``알비(Albi)''에 전송하여 실제 하드웨어를
제어한다. ``화성 탐사 로버''라는 서사를 통해 순차 실행, 반복, 조건, 변수 등
프로그래밍의 기초 개념을 단계적으로 학습하도록 설계되었다.

\begin{infobox}[title=프로젝트 핵심 특징]
\begin{itemize}
  \item \textbf{대상}: 초등학교 저\textbf{·}고학년 (블록 코딩 입문 학습자)
  \item \textbf{언어}: UI 텍스트, 코드 주석, 문서 모두 한국어
  \item \textbf{제어 대상}: 라즈베리파이 피코(RP2040) + BLE UART 모듈(HM-10/BT05)
  \item \textbf{배포 방식}: 웹(HTTPS) 호스팅 및 오프라인 \code{file://} 단독 실행 모두 지원
\end{itemize}
\end{infobox}

\subsection{세 가지 런타임 도메인}

ARES는 역할이 명확히 분리된 세 개의 실행 영역으로 구성된다.

\begin{enumerate}
  \item \textbf{웹 UI (\code{Web/})} --- 정적 HTML/JS로 작성된 브라우저 애플리케이션.
        Google Blockly로 블록 코딩 편집기를 제공하고, Web Bluetooth\footnote{\textbf{Web
        Bluetooth}: 웹 브라우저에서 자바스크립트로 BLE 기기와 직접 연결\textbf{·}통신할 수
        있게 하는 표준 웹 API. 보안상 HTTPS 또는 \code{localhost} 등 보안 컨텍스트에서만
        동작한다.} API로 로버와 통신한다. ES 모듈\footnote{\textbf{ES 모듈(ECMAScript
        Module)}: 자바스크립트의 표준 모듈 시스템. \code{import}/\code{export}로 코드를 파일
        단위로 분리하고 서로 가져다 쓸 수 있게 하여, 큰 프로그램을 책임별로 나눈다.}(\code{import}/\code{export}) 구조이며 진입점은
        \code{main.html}(블록 편집기), \code{dashboard.html}(시스템 모니터링),
        \code{index.html}(3D 랜딩 페이지)이다.
  \item \textbf{피코 펌웨어 (\code{Pico/})} --- 라즈베리파이 피코에서 동작하는
        MicroPython\footnote{\textbf{MicroPython}: 마이크로컨트롤러처럼 메모리\textbf{·}성능이
        제약된 환경에서 동작하도록 만든 경량 파이썬 구현. 표준 파이썬 문법의 부분집합을 지원한다.}
        코드. \code{main.py}가 UART\footnote{\textbf{UART(Universal Asynchronous
        Receiver/Transmitter)}: 두 장치가 별도의 클록 신호 없이 약속된 속도(보율, baud)로
        한 바이트씩 직렬로 주고받는 비동기 시리얼 통신 방식.}를 통해 BLE\footnote{\textbf{BLE(Bluetooth
        Low Energy)}: 저전력 블루투스. 적은 전력으로 소량의 데이터를 주고받도록 설계된 근거리
        무선 통신 규격.} 명령을 수신하면
        \code{CommandProcessor}(\code{process\_data.py})가 이를 해석하여
        하드웨어 모듈로 분배한다. 모든 하드웨어는 싱글턴\footnote{\textbf{싱글턴(Singleton)}:
        클래스의 인스턴스를 프로그램 전체에서 단 하나만 생성해 공유하도록 강제하는 설계 패턴.
        하드웨어처럼 실체가 하나뿐인 자원을 다룰 때 쓴다.} \code{robot}
        (\code{hardware.py})을 통해 접근한다.
  \item \textbf{AI 모듈 (\code{AI/}, \code{Web/ai\_helper*.js})} --- 자연어를
        블록으로 변환하는 보조 기능. 로컬 LLM\footnote{\textbf{LLM(Large Language Model,
        대규모 언어 모델)}: 방대한 텍스트로 학습되어 자연어 질문에 답하거나 문장을 생성하는
        인공지능 모델.} 기반 PoC\footnote{\textbf{PoC(Proof of Concept, 개념 증명)}:
        아이디어의 실현 가능성을 보이기 위한 시험적\textbf{·}초기 구현.}(\code{AI/ai.py})와,
        브라우저에서 동작하는 문법 기반 자연어→블록 파서
        (\code{Web/ai\_helper.js}, \code{ai\_helper\_grammar.js})로 구성된다.
\end{enumerate}

\subsection{데이터 흐름 개요}

전체 시스템의 명령 흐름은 다음과 같다. 웹 UI에서 생성된 텍스트 명령은
Web Bluetooth를 거쳐 피코의 UART로 전달되고, 펌웨어가 이를 해석하여
하드웨어를 구동한 뒤 응답을 역방향으로 반환한다.

\begin{lstlisting}[style=json, caption={명령 데이터 흐름}]
[웹 UI: Blockly 블록]
      | commandexecutor.js  ->  명령 문자열 생성 (예: "SERVO_tFORWARD,2")
      v
[bluetooth.js: Web Bluetooth]
      | UART 서비스(0xFFE0) / 20바이트 청크 / 개행 종단
      v
[HM-10 / BT05 BLE 모듈]  --(UART 9600 baud)-->  [피코]
      v
[main.py: UART 수신 루프]  ->  버퍼링(개행까지)  ->  CommandProcessor
      v
[process_data.py]  ->  하드웨어 모듈(wheel/leds/buzzer/sensor/oled ...)
      v
응답(예: "DIST:23.5") ->  BLE 20바이트 청크 ->  웹 UI 로그/변수 갱신
\end{lstlisting}

이 흐름을 도식화하면 그림~\ref{fig:overview}와 같다. 세 도메인은 서로 다른
언어(JavaScript, MicroPython)와 실행 환경에서 동작하지만, \emph{개행으로
구분되는 텍스트 명령 프로토콜}이라는 단일 계약으로 느슨하게 결합된다.
이러한 ``공유 코드 없는 텍스트 계약'' 구조는 양측을 독립적으로 개발\textbf{·}교체할
수 있게 하는 동시에, 한쪽의 프로토콜 변경이 반드시 다른 쪽에 반영되어야 하는
의존을 만든다(상세는 9장 참조).

\begin{diagram}{ARES 시스템 개요 및 데이터 흐름 연계도}{fig:overview}
\node[dom, text width=28mm] (web) {웹 UI\\{\footnotesize\mdseries Blockly · Web Bluetooth}};
\node[boxw, right=14mm of web] (ble) {BLE 모듈\\{\scriptsize HM-10 / BT05}};
\node[dom, text width=28mm, right=14mm of ble] (pico) {피코 펌웨어\\{\footnotesize\mdseries CommandProcessor}};
\node[hw, below=12mm of pico, text width=30mm] (hw) {하드웨어\\{\scriptsize 모터·센서·LED·부저·OLED}};

\draw[flow] (web.north east) to[bend left=18] node[lbl, above]{명령(텍스트)} (ble.north west);
\draw[flow] (ble.north east) to[bend left=18] node[lbl, above]{UART 9600} (pico.north west);
\draw[back] (ble.south west) to[bend left=18] node[lbl, below]{20B 청크 응답} (web.south east);
\draw[back] (pico.south west) to[bend left=18] (ble.south east);
\draw[flow] (pico) -- node[lbl, right]{GPIO/PWM/I2C} (hw);
\draw[back] (hw.west) to[bend left=20] node[lbl, left]{센서값} (pico.south west);

\begin{scope}[on background layer]
  \node[grp, fit=(pico)(hw), label={[font=\scriptsize\sffamily\color{rulegray}]below:로버(피코) 측}] {};
\end{scope}
\end{diagram}

\subsection{저장소 최상위 구성}

\begin{longtable}{L{3.2cm} L{11.3cm}}
\toprule
\rowcolor{tablehead}\textbf{경로} & \textbf{역할} \\
\midrule
\endhead
\code{Web/} & Blockly 블록 편집기, Web Bluetooth 통신, WebGL 시뮬레이터, 정적 자산 \\
\code{Pico/} & MicroPython 펌웨어(부팅·명령 처리·하드웨어 드라이버·설정) \\
\code{AI/} & 로컬 LLM 기반 코드 도우미 PoC (\code{ai.py}) \\
\code{Document/} & API 문서, 블록 가이드, 모듈 관리, 빌드 가이드, 본 분석 보고서 \\
\code{Missions/} & 12차시 미션 교재 LaTeX 소스 및 PDF \\
\code{Papers/} & 연구 논문 \textbf{·} 발표 자료 \\
\code{KSWeb/} & 배포/아카이브용 웹 산출물 \\
\code{XML Presets/} & 예제 Blockly 프로그램(XML) \\
\code{build.ps1 / .sh / .bat} & 오프라인 단독 실행본 빌드 자동화 스크립트 \\
\bottomrule
\end{longtable}

\subsection{보고서 구성}

본 보고서는 다음 순서로 전개된다. 2장에서 하드웨어·소프트웨어 구성요소를
조감하고, 3\textasciitilde 4장에서 하드웨어 핀 배치와 피코 펌웨어를 코드
수준에서 분석한다. 5\textasciitilde 8장은 웹 서비스 계층, Web Bluetooth,
블록 코딩 구현, WebGL 디지털 플랫폼을 다루며, 각 장은 ``파일 기반 웹앱
제작 시 요구사항''을 함께 정리한다. 9장은 코드 간 연계 구조를, 10장은
네이티브 앱 패키징 방안을, 11\textasciitilde 12장은 개선 계획과 기술적
파급효과 및 향후 개발 계획을 제시한다.

\begin{teachbox}{프로젝트 전체 그림 읽기}
\teachhead{설계 의도}
\begin{itemize}
  \item ARES는 웹\textbf{·}펌웨어\textbf{·}AI 세 영역을 \emph{텍스트 명령}이라는 단순한 약속만으로 연결한다. 서로의 내부 구현을 몰라도 되게 하여 독립적 개발\textbf{·}교체가 쉽다.
  \item 인터넷\textbf{·}설치가 어려운 교실을 전제로 \emph{오프라인 우선}으로 설계했다.
\end{itemize}
\teachhead{분석할 때 핵심 질문}
\begin{itemize}
  \item 블록 하나를 누르면 명령은 \emph{어디서} 만들어져 \emph{어떤 경로로} 로버에 도달하는가?
  \item 웹과 펌웨어가 코드를 공유하지 않는데 어떻게 서로 통하는가?
\end{itemize}
\teachhead{점검 요소}
\begin{itemize}
  \item 좋아하는 블록 1개를 골라 ``블록 $\to$ 명령 문자열 $\to$ BLE $\to$ 펌웨어 핸들러 $\to$ 하드웨어''의 전 경로를 직접 그려 보고, 2장 이후 내용과 대조하라.
\end{itemize}
\end{teachbox}

